Parallel exchange rate markets are a defining feature of economies operating under fixed or heavily managed exchange rate regimes. When the official exchange rate becomes misaligned with underlying fundamentals, excess demand for foreign currency emerges and is rationed administratively. This gives rise to a parallel market in which foreign currency trades at a premium relative to the official rate.

Traditionally, these markets have relied on physical foreign currency and informal intermediaries. However, recent years have witnessed the rapid emergence of stablecoins as an alternative vehicle for accessing foreign currency. Stablecoins are digital tokens pegged to major currencies and traded on decentralized platforms. They offer lower transaction costs, near-instant settlement, and global accessibility.
